\section*{\centering TÓM TẮT}

\textbf{Tóm tắt:} Bài toán cái túi hợp tập (SUKP) là một biến thể NP-khó của bài toán cái túi, trong đó mỗi item là một tập element có trọng số và trọng số của element dùng chung chỉ được tính một lần. Các thuật toán tìm kiếm cục bộ tốt nhất cho SUKP, như E2LS, dựa vào một luật viết tay là tỉ số giá trị trên trọng số để quyết định thêm item nào. Luận văn này giữ nguyên khung E2LS và chỉ thay đúng một thành phần, bước xếp hạng ứng viên của toán tử ADD*, bằng một hàm được học. Luận văn cài đặt lại E2LS bằng C++ với giao diện Python cho phép thay hàm xếp hạng, thiết kế một mạng nơ-ron đồ thị trên đồ thị hai phía item--element với số tham số không phụ thuộc kích thước instance, và huấn luyện hàm xếp hạng theo hai cách: học tăng cường bằng DQN trên từng instance, và học bắt chước từ nhãn tối ưu được tính bằng tìm kiếm vét cạn trên các instance sinh ngẫu nhiên. Bộ xếp hạng học bắt chước được tham số hóa như một hiệu chỉnh của luật tỉ số, nên khi chưa huấn luyện nó trùng với thuật toán gốc. Trên các instance kiểm định chưa gặp, bộ xếp hạng nâng tỉ lệ chọn đúng ứng viên tối ưu từ 60.9\% lên 71.3\% và giảm regret từ 0.71\% xuống 0.44\%. Với một lần chạy 30 giây, bản cài đặt E2LS đạt lời giải tốt nhất đã biết trên 17/30 instance của Set I và có khoảng cách trung bình 1.45\% trên Set II. Luận văn cũng phân tích chi phí suy luận của việc xếp hạng bằng mạng tại mỗi nút tìm kiếm và đề xuất quy trình đánh giá cùng số bước và cùng thời gian cho các hàm xếp hạng học được.

\textbf{Từ khóa:} bài toán cái túi hợp tập, tìm kiếm cục bộ, học tăng cường, mạng nơ-ron đồ thị, học máy cho tối ưu tổ hợp.

\vspace{0.8cm}

\section*{\centering ABSTRACT}

\textbf{Abstract:} The Set-Union Knapsack Problem (SUKP) is an NP-hard variant of the knapsack problem in which each item is a set of weighted elements and a shared element is paid only once. The best local search algorithms for SUKP, such as E2LS, rely on a hand-written value-weight ratio to decide which item to add. This thesis keeps the E2LS framework unchanged and replaces exactly one component, the candidate ranking inside the ADD* operator, with a learned function. We reimplement E2LS in C++ with a Python interface that accepts any ranking function, design a graph neural network over the item--element bipartite graph whose parameter count does not depend on the instance size, and train the ranking function in two ways: per-instance deep Q-learning, and imitation of optimal labels computed by exhaustive search on randomly generated instances. The imitation ranker is parameterised as a correction to the ratio rule, so an untrained model reproduces the original algorithm. On unseen validation instances, the learned ranker raises the rate of picking an optimal candidate from 60.9\% to 71.3\% and lowers the regret from 0.71\% to 0.44\%. With a single 30-second run, the reimplemented E2LS matches the best-known solution on 17 of the 30 Set I instances and is within 1.45\% of it on average on Set II. We also analyse the inference cost of ranking with the network at every search node, and propose an equal-steps and equal-time protocol for evaluating learned rankers on the benchmark.

\textbf{Keywords:} set-union knapsack problem, local search, reinforcement learning, graph neural network, machine learning for combinatorial optimization.

\newpage
